\documentclass[french]{article}

%Chargement des paquets
\usepackage{amsmath}
\usepackage{amsfonts}
\usepackage{amssymb}
\usepackage{enumerate}
\usepackage{mathtools}
\usepackage{bbm}
\usepackage{xparse, etoolbox}
\usepackage{enumerate}
\usepackage{enumitem}
\usepackage{mathabx}
\usepackage{babel}[french]

%environment
\usepackage{amsthm}
\usepackage{keytheorems}
\usepackage{hyperref} 

%Interface théorème
\newkeytheorem{lem}[
	name = Lemme
]

\newkeytheorem{prop}[
	name = Proposition
]

\newkeytheorem{thm}[
	name = Théorème
]

\newkeytheorem{coro}[
	name = Corollaire
]

\renewcommand*{\proofname}{Démonstration}

\theoremstyle{definition}
\newtheorem{defn}{Définition}

\theoremstyle{definition}
\newtheorem*{nota}{Notation}

\theoremstyle{definition}
\newtheorem*{limi}{Liminaire}

\theoremstyle{remark}
\newtheorem{rmq}{Remarque}

\theoremstyle{plain}
\newtheorem*{fait}{Fait}


%Ensembles de nombres
\newcommand{\N} {\mathbb{N}}
\newcommand{\Ne}{\N^\ast}
\newcommand{\Z} {\mathbb{Z}}
\newcommand{\Q} {\mathbb{Q}}
\newcommand{\R} {\mathbb{R}}
\newcommand{\Rb}{\overline{\mathbb{R}}}
\newcommand{\Rp}{\R_+}
\newcommand{\Rm}{\R_-}
\newcommand{\K} {\mathbb{K}}
\newcommand{\Ket} {\mathbb{K^{\ast}}}
\newcommand{\Cx}{\mathbb{C}}

%Ensembles, tribus
\newcommand{\A}{\mathbb{A}}
\renewcommand{\P}{\mathcal{P}}
\newcommand{\T}{\mathcal{T}}
\newcommand{\F}{\mathcal{F}}
\newcommand{\C} {\mathcal{C}}
\newcommand{\ouv}{\mathcal{O}}
\newcommand{\B}{\mathcal{B}}
\newcommand{\M}{\mathcal{M}}
\newcommand{\etag}{\mathcal{E}}
\newcommand{\etagOp}{\etag^{+}(\Omega)}
\newcommand{\etagO}{\etag(\Omega)}
%%Lp avant quotientage
\newcommand{\Lic}{\mathcal{L}^1}
\newcommand{\Lpc}{\mathcal{L}^p}
\newcommand{\Linfc}{\mathcal{L}^{\infty}}
\newcommand{\Lifull}{\Lic(\Omega, \B, m)}
\newcommand{\Lpfull}{\Lpc(\Omega, \B, m)}
\newcommand{\Linffull}{\Linfc(\Omega, \B, m)}
%%Lp après quotientage
\newcommand{\Li}{L^1}
\newcommand{\Ld}{L^2}
\newcommand{\Lp}{L^p}
\newcommand{\Linf}{L^{\infty}}
\newcommand{\Listd}{\Li(\Omega, m)} 
\newcommand{\Ldstd}{\Ld(\Omega, m)} 
\newcommand{\Lpstd}{\Lp(\Omega, m)} 

%Quelques opérateurs
\DeclareMathOperator{\codim}{codim}
\DeclareMathOperator{\rg}{rg}
\DeclareMathOperator{\tr}{tr}
\DeclareMathOperator{\vect}{Vect}
\DeclareMathOperator{\ima}{Im}
\DeclareMathOperator{\id}{Id}
\DeclareMathOperator{\sign}{sign}
\DeclareMathOperator{\ext}{ext}
\newcommand{\ind}{\mathbbm{1}}
\newcommand*{\spr}[2]{\left(\,#1\,\middle|\,#2\,\right)}
\newcommand{\interior}[1]{%
  {\kern0pt#1}^{\mathrm{o}}%
}
\DeclareMathOperator{\card}{card}
\DeclareMathOperator{\ppcm}{ppcm}

%Commandes de pure flemme
\newcommand{\veps}{\varepsilon}
\newcommand{\intab}{\int_{a}^{b}}
\newcommand{\intR}{\int_{-\infty}^{\infty}}
\newcommand{\intinf}{\int_{- \infty}^{+ \infty}}
\newcommand{\equi}{\Leftrightarrow}
\newcommand{\Kev}{$\K$-espace vectoriel }
\newcommand{\Kevs}{$\K$-espaces vectoriels }
\newcommand{\Rev}{$\R$-espace vectoriel }
\newcommand{\Revs}{$\R$-espaces vectoriels }
\newcommand{\Cev}{$\Cx$-espace vectoriel }
\newcommand{\Cevs}{$\Cx$-espaces vectoriels }
\newcommand{\evn}{(E, \norm{.})}
\newcommand{\interf}[2]{[#1,#2]}
\newcommand{\limb}{\lim\limits_}
\newcommand{\lebn}{\lambda_n}
\newcommand{\pp}{\text{~presque partout}}
\newcommand{\derivp}[2]{\frac{\partial #1}{\partial #2}}
\newcommand{\famiu}{(u_i)_{i \in I}}
\newcommand{\sev}{\text{~s.e.v.}}
\newcommand{\Mnp}{M_{n,p}(\K)}
\newcommand{\Mn}{M_{n}(\K)}
\newcommand{\tq}{\text{~t.q.~}}
\newcommand{\mq}{~montrer~que~}
\newcommand{\et}{\quad \text{et} \quad}
\newcommand{\ou}{\text{~ou~}}
\newcommand{\Kx}{\K[X]}


%Commandes de gars chelou
\renewcommand{\subset}{\subseteq}

%Commande restriction, ty duckduckgo
\def\restr#1#2{\mathchoice
              {\setbox1\hbox{${\displaystyle #1}_{\scriptstyle #2}$}
              \restraux{#1}{#2}}
              {\setbox1\hbox{${\textstyle #1}_{\scriptstyle #2}$}
              \restraux{#1}{#2}}
              {\setbox1\hbox{${\scriptstyle #1}_{\scriptscriptstyle #2}$}
              \restraux{#1}{#2}}
              {\setbox1\hbox{${\scriptscriptstyle #1}_{\scriptscriptstyle #2}$}
              \restraux{#1}{#2}}}
\def\restraux#1#2{{#1\,\smash{\vrule height .8\ht1 depth .85\dp1}}_{\,#2}} 

%Quelques normes
\DeclarePairedDelimiter\abs{\lvert}{\rvert}
\DeclarePairedDelimiter\labs{\bigl\lvert}{\bigr\rvert}
\DeclarePairedDelimiter\norm{\lVert}{\rVert}
\DeclarePairedDelimiterXPP\normi[1]{}\lVert\rVert{_1}{\ifblank{#1}{\:\cdot\:}{#1}}
\DeclarePairedDelimiterXPP\normd[1]{}\lVert\rVert{_2}{\ifblank{#1}{\:\cdot\:}{#1}}
\DeclarePairedDelimiterXPP\normp[1]{}\lVert\rVert{_p}{\ifblank{#1}{\:\cdot\:}{#1}}
\DeclarePairedDelimiterXPP\normq[1]{}\lVert\rVert{_q}{\ifblank{#1}{\:\cdot\:}{#1}}
\DeclarePairedDelimiterXPP\norminf[1]{}\lVert\rVert{_\infty}{\ifblank{#1}{\:\cdot\:}{#1}}

%Bracket en angle
\DeclarePairedDelimiter\angl{\langle}{\rangle}

%Set, parenthèses
\newcommand{\set}[1]{\{ #1 \}}
\newcommand{\bigset}[1]{\bigl\{ #1 \bigr\}}
\newcommand{\Bigset}[1]{\Bigl\{ #1 \Bigr\}}
\newcommand{\bigpar}[1]{\bigl( #1 \bigr)}
\newcommand{\Bigpar}[1]{\Bigl( #1 \Bigr)}

%Opitmisation des opérateurs de comparaison
\DeclarePairedDelimiter\tio{o (}{)}
\DeclarePairedDelimiter\gro{O (}{)}


\begin{document}

\begin{defn}[Série entière]
Une série entière est une série numérique de la forme $\sum a_n z^n$, où $(a_n)_\N$ est une suite de nombres complexes.
En désignant par D l’ensemble des nombres complexes z pour lesquels la série $z \mapsto a_n z^n$ est convergente, on note,
pour tout $z \in D, S(z) =  \sum_{n=0}^{\infty} a_n z^n$ la somme de cette série, et on définit ainsi la fonction $S$ de D dans $\Cx$.
\end{defn}

\section{Rayon de convergence}

\begin{thm}[Abel]
Soit $\sum a_n z^n$ une série entière. S'il existe un scalaire non nul $z_0$ tel que la suite $(a_n z_0^n)_\N$ soit bornée, alors la 
série entière $\sum a_n z^n$ converge absolument pour tout nombre complexe $z$ tel que $\abs{z} < \abs{z_0}$.
\end{thm}

\begin{proof}
Pour $N \in \N$ fixé, on a :
\begin{align*}
\sum_{n=0}^N a_n \abs{z_n}^n = \sum _{n=0}^N a_n \abs{z_0}^n \dfrac{\abs{z_n}^n}{\abs{z_0}^n} \\
& \leq \sum _{n=0}^N M \dfrac{\abs{z_n}^n}{\abs{z_0}^n} ,
\end{align*}
où $M$ est un majorant de la suite $(a_n z_0^n)_\N$. Pour $\abs{z} < \abs{z_0}$, la série converge (série géométrique).
\end{proof}

\begin{thm}
Soit $\sum a_n z^n$ une série entière. L'ensemble $I$ des réels positifs ou nuls tels que la suite $(a_n r^n)$ soit bornée est un 
intervalle qui est, soit réduit à $\{0\}$, soit de la forme $[0,R]$ ou $[0,R[$ avec $R >0$, soit égal à $\Rp$ tout entier.
\end{thm}

\begin{proof}
Il est évident que la série $\sum a_n z^n$ converge pour $z=0$, donc $I$ contient toujours au moins $0$. \\
Supposons que $I \neq \set{0}$, selon le théorème précédent, si $r \in I$ alors $I$ contient $[0,r]$ tout entier. En effet, pour tout 
réel $s \in [0,r[$, la série $\sum a_n s^n$ converge donc la suite $(a_n s^n)_\N$ converge vers 0, ce qui implique qu'elle est bornée. \\
Finalement, l'ensemble $I$ est un intervalle de $\Rp$ qui contient 0, il a donc une des formes proposées.
\end{proof}

\begin{rmq}
En particulier, $I$ admet une borne supérieure.
\end{rmq}

\begin{defn}
La rayon de convergence de la série entière $\sum a_n z^n$ est l'élément de $\widebar{\Rp}$ défini par :
\[ R = \sup \set{ r \in \R+ \tq (a_n r^n)_\N \text{ est bornée} } . \]
\end{defn}

\begin{rmq}
Si $\abs{z} > R$, alors $\abs{a_n}\abs{z}^n$ ne tend pas vers 0 donc la série $\sum a_n z^n$ diverge trivialement.
\end{rmq}

\begin{thm}
Soient $\sum a_n z^n$ et $\sum b_n z^n$ deux séries entières de rayons de convergence respectifs $R$ et $R'$.
\begin{enumerate}[label=(\roman)]
\item Si $\abs{a_n} \leq \abs{b_n}$ pour tout $n \in \N$, alors $R \geq R'$.
\item Si $a_n = \gro{b_n}$, alors $R \geq R'$.
\item Si $a_n \sim b_n$, alors $R = R'$.
\end{enumerate}
\end{thm}

\begin{proof}
\begin{enumerate}[label=(\roman)]
\item Clairement :
\[ \set{ r \in \R+ \tq (b_n r^n)_\N \text{ est bornée} } \subset \set { r \in \R+ \tq (a_n r^n)_\N \text{ est bornée} } \]
d'où l'inégalité sur les $\sup$.  
\item Si $a_n = \gro{b_n}$, alors il existe une suite $(\varpi(n))_\N$ bornée tel que $a_n = \varphi_n b_n$ à patir d'un certain rang.
 Donc, en notant $M >0$ le réel tel que $\forall n \in \N, \abs{\varphi(n)} \leq M$ on déduit que 
$\abs{a_n} \leq M \abs{b_n}$ pour tout $n \in \N$. Or
 \[ \set{ r \in \R+ \tq (b_n r^n)_\N \text{ est bornée} } = \set{ r \in \R+ \tq (M b_n r^n)_\N \text{ est bornée} } , \]
 d'où le résultat, selon (i).
\item Si $a_n \sim b_n$, alors il existe une suite $(\varpi(n))_\N$ de limite égale à 1 telle que $a_n = \varphi_n b_n$ à partir d'un 
certain rang. Ce qui assure que $a_n = \gro{b_n}$ et $b_n = \gro{a_n}$ à partir d'un certain rang (afin d'éviter toute annulation de 
$\varphi_n$). D'où le résultat selon (ii). 
\end{enumerate}
\end{proof}

\begin{thm}
Soit $\sum a_n z^n$ une série entière de rayon de convergence $R$.
\begin{enumerate}[label=(\roman)]
\item si $R > 0$, la série $\sum a_n z^n$ est absolument convergente pour tout $z$ tel que $\abs{z} < R$ ;
\item si $R$ est fini, la série $\sum a_n z^n$ est divergente pour tout $z$ tel que $\abs{z} > R$.
\end{enumerate}
\end{thm}

\begin{proof}
\begin{enumerate}[label=(\roman)]
\item Si $R>0$ et $\abs{z} < R$, alors il existe $r \in \R$ tel que $z < r < R$ et $r$ est dans l'ensemble $I$ défini plus tôt.
Donc, d'après le théorème d'Abel, la série $\sum a_n z^n$ est absolument convergente.
\item si $R$ est fini et $\abs{z} > R$, alors $\abs{z} \not \in I$ donc la suite $(a_n z^n)$ n'est pas bornée, par suite, la série diverge
trivialement.
\end{enumerate}
\end{proof}

\begin{thm}[D'Alembert]
Soit $\sum a_n z^n$ une série entière telle que $a_n \neq 0$ à partir d'un certain rang. Si
\[ \lim \abs*{\dfrac{a_{n+1}}{a_n}} = l \in [0, \infty], \]
alors le rayon de convergence de la série est $\dfrac1{l}$ (avec $1/0 = \infty$ et $1/\infty = 0$). 
\end{thm}

\begin{proof}
Il s'agit d'une application directe de la règle de D'Alembert pour les séries, en effet :
\[ \lim \abs*{\dfrac{a_{n+1}z^{n+1}}{a_nz^n}} = \abs{z} l . \]
Ainsi, si $\abs{z} < \frac{1}{l}$ la série converge. De même si $\abs{z} > \frac{1}{l}$, alors la série diverge.
\end{proof}

\begin{thm}[Cauchy]
Soit $\sum a_n z^n$ une série entière. Si 
\[ \lim \abs*{a_n}^{1/n} = l \in [0, \infty], \]
alors le rayon de convergence de la série est $\dfrac1{l}$. 
\end{thm}

\begin{proof}
Il s'agit à nouveau d'une application du critère de Cauchy à la série de terme général $a_n z^n$.
\end{proof}

\begin{thm}[Hadamard]
Soit $\sum a_n z^n$ une série entière. Si 
\[ \limsup \abs*{a_n}^{1/n} = l \in [0, \infty], \]
alors le rayon de convergence de la série est $\dfrac1{l}$.
\end{thm}

\begin{proof}
Supposons dans un premier temps que $l \in ]0, \infty[$ et soit $z \in \Cx$ tel que $0 < \abs{z} < \frac1{l}$.
Il existe donc $b \in \Rp$ tel que $0 < l < b < \frac{1}{\abs{z}}$.
Comme $l = \limsup \abs*{a_n}^{1/n}$, il existe $N \in \N$ tel que $\forall n \geq N, \abs*{a_n}^{1/n} < b$.
Autrement dit 
\[ \exists N \in \N, \forall n \in \N : n \geq N \Rightarrow \abs*{a_n} \abs{z}^n < b^n \abs{z}^n = (b\abs{z})^n . \]
Or, $b\abs{z} < 1$, donc $\sum a_n z^n$ converge absolument. \\
Pour $z \in \Cx$ tel que $\abs{z} > \frac1{l}$, on procède de la même manière. Il existe $b \in \R$ tel que 
$0 < \frac{1}{\abs{z}} < b < l$ et 
\[ \forall N \in \N, \exists n \in \N \tq n \geq N \text{ et } \abs*{a_n} \abs{z}^n \geq b^n \abs{z}^n = (b\abs{z})^n . \]
Or, $b\abs{z} > 1$, donc la suite $(\abs{a_n} \abs{z}^n)$ ne tend pas vers 0 et la série diverge trivialement. \newline

Si $l=0$, alors pour tout $z \in \Cx$ non nul, il existe $b \in \Rp$ tel que $0 < b < \frac{1}{\abs{z}}$. 
Par le même raisonnement, on conclut que la série converge. \newline

Si $l=\infty$, considérons $z \in \Cx$ non nul. On a :
\[ \forall  N \in \N, \exists n \in \N \tq n \geq N \text{ et } \abs*{a_n}^{1/n} > \dfrac1{\abs{z}} . \]
Donc la suite $(\abs{a_n} \abs{z}^n)$ ne tend pas vers 0 et la série diverge trivialement. 
Donc le rayon de convergence est strictement inférieur à $\abs{z}$ pour tout $z$ non nul, donc $R=0$.
\end{proof}

\section{Opérations sur les séries entières}

\begin{defn}[Série somme]
La somme des séries entières $\sum a_n z^n$ et $\sum b_n z^n$ est $\sum (a_n + b_n) z^n$
\end{defn}

\begin{thm}
Soient $\sum a_n z^n$, $\sum b_n z^n$ deux séries entières de rayons de convergence respectifs $R$ et $R'$. On note aussi
$R''$ le rayon de convergence de la série entière somme $\sum (a_n + b_n) z^n$.
Pour $R \neq R'$ on a $R'' = \min(R, R')$ et pour $R=R'$ on a $R'' \geq \min(R,R')$. 
Dans tous les cas, pour $\abs{z} < \min(R, R')$ on a :
\[ \sum_{n=0}^{\infty}(a_n+b_n) z^n = \sum_{n=0}^{\infty}a_n z^n + \sum_{n=0}^{\infty}b_n z^n . \]
\end{thm}

\begin{proof}
Il est clair que si $\abs{z} < \min(R, R')$ la série somme est absolument convergente comme somme de deux séries absolument 
convergente et on a bien la formule annoncée. On a donc $R'' \geq \min(R,R')$. \\
Supposons que $R \neq R'$, par exemple $0 \leq R < R'$. On peut alors trouver un réel $r$ tel que $R < r < R'$ et la série 
$\sum (a_n + b_n)r^n$ est divergente comme somme de la série divergente $\sum a_n r^n$ et de la série convergente 
$\sum r_n z^n$. Donc $R'' \leq R$ soit $R'' = R$ dans ce cas.
\end{proof}

\begin{defn}[Série produit]
Soient $\sum a_n z^n$ et $\sum b_n z^n$ des séries entières, la série entière produit est la série $\sum c_n z^n$ avec 
\[ c_n = \sum_{k=0}^n a_k b_{n-k} . \]
\end{defn}

\begin{thm}
Soit $\sum a_n z^n$ une série entière de rayon de convergence $R$. La série dérivée $\sum n a_n z^{n-1}$ et la série primitive 
$\sum \dfrac{a_n}{n+1} z^{n+1}$ ont le même rayon de convergence égal à $R$.
\end{thm}

\begin{proof}
Il suffit d'étudier le rayon de convergence de la sériee dérivée (car la dérivée de la série primitive est la série entière). On note :
\[ I = \set{ r \in \R+ \tq (a_n r^n)_\N \text{ est bornée} } \et I' = \set{ r \in \R+ \tq (n a_n r^{n-1})_\N \text{ est bornée} } . \]
Ce sont des intervalles de $\Rp$ contenant 0 et de bornes supérieures respectives $R$ et $R'$.
Or, pour $n \geq 1$, on a $\abs{a_n r^n} \leq \abs{n a_n r^n} = r \abs{n a_n r^{n-1}$, on en déduit que $R' \leq R$.
Si $R=0$, alors $R'=0$, sinon pour $r \in ]0, R[$, on peut trouver $s$ tel que $r < s < R$, et en posant $h = s - r >0$, on a :
\[ \forall n \geq 1, s^n = (r+h)^n = \sum_{k=0}^n \binom{k}{n} r^{n-k} h^k \geq n r^{n-1} h , \]
et donc $\abs{n a_n r^{n-1} \leq \dfrac1{h} \abs{a_n s^n}$. Or $s \in I$, donc la suite $(a_n s^n)_\N$ est bornée et en conséquence,
$r \in I'$. Ceci étant vrai pour tout $r \in ]0, R[$, on a donc $]0, R[ \subset I'$ ce qui entraîne $R \leq R'$ et $R=R'$.
\end{proof}

\section{Fonctions développables en série entière}

\begin{defn}

\end{defn}

%\begin{rmq}
%A noter que si une série $S = \sum a_n z^n$ a pour rayon de convergence $R$, alors $\limsup \abs*{a_n}^{1/n}=1/R$.
%\end{rmq}
%
%\section{Propriétés analytiques}
%
%\begin{lem}[Abel]
%Soit $S = \sum a_n z^n$ où $(a_n)_\N \subset \Cx$, une série entière. On suppose que son rayon de convergence $R \in [0, \infty]$
%est strictement positif. Pour tout $A \in ]0, R[$, la série converge uniformément sur le disque fermé de centre $O$ et de rayon $A$ :
%$D(0, A) = \set{z \in \Cx ; \abs{z} \leq A}$.
%\end{lem}
%
%\begin{proof}
%On va démontrer que la convergence est même normale sur le disque, en effet :
%\[ \sup_{D(0,A)} \abs{a_n z^n} = \abs{a_n} A^n , \]
%et comme $A<R$, $\sum \abs{a_n}A^n$ converge.
%\end{proof}
%
%\begin{defn}
%Soit $S = \sum a_n z^n$ une série entière. On appelle série dérivée de $S$, et on note $S'$, la série définie par :
%\[ S' = \sum (n+1) a_{n+1} z^n . \]
%\end{defn}
%
%\begin{thm}
%\label{thm:rayonderivée}
%La série dérivée de $S$ a le même rayon de convergence que $S$.
%\end{thm}
%
%\begin{proof}
%On se ramène à l'étude de la série $\sum n a_n z^n$ en multipliant $S'$ par $z$ (ce qui ne change pas son rayon de convergence). \\
%On remarque que $\lim n^{1/n} = \lim \exp(\frac{1}{n} \ln(n)) = 1$. \\
%Donc, pour tout $\veps \in ]0, 1[$, il existe $N \in \N$ tel que $\forall n \geq N, n^{1/n} \in [1-\veps, 1+\veps]$. \\
%A partir du rang $N$, on peut donc écrire :
%\[ (1-\veps) \abs{a_{n}}^{1/n} \leq \abs{na_{n}}^{1/n} \leq (1+\veps) \abs{a_{n}}^{1/n} \]
%soit, en passant à la limite supérieure :
%\[ (1-\veps) \limsup \abs{a_{n}}^{1/n} \leq \limsup \abs{na_{n}}^{1/n} \leq (1+\veps) \limsup \abs{a_{n}}^{1/n} . \]
%Ceci étant vrai pour tout $\veps \in ]0,1[$, on conclut que :
%\[ \limsup \abs{a_{n}}^{1/n} = \limsup \abs{na_{n}}^{1/n} \]
%ce qui achève la démonstration selon le critère de Hadamard.
%\end{proof}
%
%\begin{coro}
%Soit $S : x \mapsto \sum_{n=0}^{\infty} a_n x^n$ la somme d'une série entière réelle de rayon de convergence $R$ strictement positif.
%Alors $S$ est de classe $\C^{\infty}$ sur $]-R, R[$. De plus, $\forall x \in ]-R, R[, \forall k \in \Ne$ :
%\[ S^{(k)}(x) = \sum_{n=k}^{\infty} \dfrac{n!}{k!} a_n x^{n-k} . \]
%\end{coro}
%
%\begin{proof}
%La suite de fonctions $\Bigpar{f_n = \sum_{k=0}^n a_k x^k}_{n \in \N}$ est de classe $\C^{\infty}$ sur $]-R, R[$. De plus :
%\begin{enumerate}[label=(\roman*)]
%\item il existe $x_0 \in ]-R, R[$ tel que la suite $(f_n(x_0))_\N$ converge (par exemple $x_0=0$) ;
%\item la suite $(f'_n)_\N$ converge uniformément sur tout segment $[-A, A] \subset ]-R, R[$ vers la série dérivée de $S$ 
%(selon le lemme d'Abel et le théorème $\ref{thm:rayonderivée}$).
%\end{enumerate}
%Donc, d'après le théorème de dérivation des suites de fonctions, la suite $(f_n)_\N$ converge uniformément vers $S$, de classe $C^1$, 
%sur tout segment  $[-A, A] \subset ]-R, R[$ donc sur $]-R, R[$. 
%En itérant ce raisonnement, on démontre que $S$ est de classe $C^{\infty}$ sur $]-R,R[$.
%\end{proof}
%
%\begin{thm}[Intégration terme à termes]
%Soit $S : x \mapsto \sum_{n=0}^{\infty} a_n x^n$ la somme d'une série entière réelle de rayon de convergence $R$ strictement positif.
%\[\forall x \in ]-R, R[, \text{ on a } \int_0^x S(t) dt = \sum_{n=0}^{\infty} \dfrac{a_n}{n+1} x^{n+1} \]
%\end{thm}
%
%\begin{proof}
%Il s'agit d'une conséquence de l'intégration terme à termes d'une suite de fonctions et du fait que les sommes partielles convergent 
%uniformément vers $S$ sur tout segment $[-A, A] \subset ]-R, R[$.
%\end{proof}
%
%\begin{thm}[Abel]
%Soit $S : x \mapsto \sum_{n=0}^{\infty} a_n x^n$ la somme d'une série entière réelle de rayon de convergence $R$ strictement positif.
%Si $S$ est définie en $R$, i.e. si $\sum_{n=0}^{\infty} a_n R^n$ converge, alors $S$ converge uniformément sur $[0,R]$.
%Et donc $S$ est continue en $R$, $\lim_{x \to \R} S(x) = S(R)$.
%\end{thm}
%
%\begin{rmq}
%On a le même résultat en $-R$ (il suffit de considérer la fonction $T : x \mapsto S(-x)$).
%\end{rmq}
%
%\begin{proof}
%On veut montrer que la série $S$ converge uniformément sur $[0,R]$, ce qui est équivalent à montrer que :
%\[ \lim_{N\to \infty} \sup_{x \in [0,R]} \abs*{\sum_{n=N}^{\infty} a_n x^n} = 0  . \]
%Comme la série numérique $S(R)$ converge, la suite $r_N = \sum_{n=N}^{\infty} a_n R^n$ tend vers 0, autrement dit :
%\[ \forall \veps > 0, \exists N_{\veps} \in \N, \forall n \in \N :~ n \geq N_{\veps} \Rightarrow \abs{r_n} \leq \veps . \]
%On va s'intéresser aux restes partiels de la série $S$ :
%\begin{align*}
%\sum_{n=p}^q a_n x^n & = \sum_{n=p}^q a_n R^n \left ( \dfrac{x}{R} \right)^n \\
%& = \sum_{n=p}^q (r_n - r_{n+1}) \left ( \dfrac{x}{R} \right )^n \text{ on utilise la transformation d'Abel} \\
%& = \sum_{n=p}^q r_n \left ( \dfrac{x}{R} \right)^n - \sum_{n=p}^q r_{n+1} \left ( \dfrac{x}{R} \right )^n \\
%& = r_p \left ( \dfrac{x}{R} \right)^p - r_{q+1} \left ( \dfrac{x}{R} \right)^q + 
%\sum_{n=p+1}^q r_n \left ( \left ( \dfrac{x}{R} \right)^n - \left ( \dfrac{x}{R} \right)^{n-1} \right) 
%\end{align*}
%Donc, pour $p \geq N_{\veps}$ :
%\begin{align*} 
%\abs*{\sum_{n=p}^q a_n x^n} & \leq \abs*{r_p \left ( \dfrac{x}{R} \right )^p } + \abs*{ r_{q+1} \left ( \dfrac{x}{R} \right )^q}  
%+ \sum_{n=p+1}^q \abs{r_n} \abs*{ \left ( \dfrac{x}{R} \right )^n - \left ( \dfrac{x}{R} \right)^{n-1} } \\
%& \leq \veps \left( \left ( \dfrac{x}{R} \right) ^p + \left ( \dfrac{x}{R} \right)^q + 
%\sum_{n=p+1}^q  \abs*{ \left ( \dfrac{x}{R} \right)^n - \left ( \dfrac{x}{R} \right)^{n-1} } \right ) \\
%& \leq \veps \left( \left ( \dfrac{x}{R} \right) ^p + \left ( \dfrac{x}{R} \right)^q + 
%\sum_{n=p+1}^q \left ( \dfrac{x}{R} \right)^{n-1} - \left ( \dfrac{x}{R} \right)^n \right ) \text{ car } \dfrac{x}{R} \leq 1 \\
%& \leq 2 \veps \left ( \dfrac{x}{R} \right) ^p \text{ somme téléscopique } \\
%& \leq 2 \veps
%\end{align*}
%On passe à la limite quand $q \to \infty$ et on conclut que :
%\[ \forall \veps > 0, \exists N_{\veps} \in \N, \forall p \in \N :~ 
%p \geq N_{\veps} \Rightarrow \forall x \in [0, R], \abs*{\sum_{n=p}^{\infty} a_n x^n} \leq 2\veps . \]
%Donc la suite converge bien uniformément sur $[0,R]$, ce qui entraîne la continuité de $S$ sur ce même intervalle. 
%\end{proof}
%
%\begin{rmq}
%A cette formulation du théorème, restreinte aux séries entières réelles, on peut préférer la formulation ci-dessous.
%\end{rmq}
%
%\begin{thm}[Abel Radial]
%Soit $\sum a_n z^n$ une série entière à coefficients complexes de rayon de convergence $R$ fini et strictement positif.
%On note $S(z)$ sa somme pour $z \in D(0, R)$.
%On suppose qu'il existe $z_0 \in S(0,R)$ tel que la série $\sum a_n z_0^n$ soit convergente. Dans ce cas, on a :
%\[ \lim{t \to 1^{-1}} \sum_{n=0}^{\infty} a_n z_0^n t^n = \sum_{n=0}^{\infty} a_n z_0^n . \]
%\end{thm}
%
%\begin{proof}
%Pour tous $n \in \Ne$ et $t \in [0,1]$, on note :
%\[ S_n(t) = \sum_{k=0}^n a_k z_0^k t^k \et \varphi(t) = S(tz_0) = \sum_{n_0}^{\infty} a_n z_0^n t^t . \]
%On note que $\varphi$ est bien définie par hypothèse. De plus :
%\begin{align*}
%S_n(t) & = a_0 + \sum_{k=1}^n a_k z_0^k t^k  \\
%& S_0(1) + \sum_{k=1}^n ( S_k(1) - S_{k-1}(1) ) t^k \text{ on utilise la transformation d'Abel} \\
%& S_0(1) + \sum_{k=1}^n S_k(1) t^k - \sum_{k=0}^{n-1} S_k(1) t^{k+1} \\
%& \sum_{k=0}^{n-1} S_k(1) ( t^k - t^{k+1} ) + S_n(1)t^n \\
%& = (1-t) \sum_{k=0}^{n-1} S_k(1) t^k + S_n(1)t^n .
%\end{align*} 
%Pour $t \in [0, 1[$, on a $\lim_{n \to \infty} S_n(1) t^n = \varphi(1) \cdot 0 = 0$, on en déduit donc que la série $\sum S_n t^n$ converge
%sur $[0,1[$ et, par passage à la limite dans l'égalité ci-dessus, que :
%\[ \varphi(t) = (1-t) \sum_{n=0}^{\infty} S_n(1) t^n , \text{ pour } t \in [0,1[ . \]
%En utilisant le fait que $\sum t^n \to \dfrac{1}{1-t}$ pour $t \in [0,1 [$ on peut écrire :
%\begin{align*}
%\varphi(t) - varphi(1) & = (1-t) \sum_{n=0}^{\infty} S_n(1) t^n - \varphi(1) (1-t) \sum_{n=0}^{\infty} t^n \\
%& = (1-t) \sum_{n=0}^{\infty} ( S_n(1) - \varphi(i)) t^n . 
%\end{align*}
%Comme $\lim S_n(1) = \varphi(1)$, pour tout réel $\veps >0$, il existe un entier $N_{\veps}$ tel que 
%\[ \forall n > N_{\veps}, \abs{S_n(1) - \varphi(1) \leq \veps , \]
%ce qui donne :
%\begin{align*}
%\abs{\varphi(t) - varphi(1)} & \leq  (1-t) \left ( \sum_{n=0}^{N_{\veps}} \abs{ S_n(1) - \varphi(i) } t^n 
%+ \veps \sum_{n=N_{\veps}}^{\infty} t^n \right ) \\
%& \leq (1-t) \sum_{n=0}^{N_{\veps}} \abs{ S_n(1) - \varphi(i) } + \veps (1-t) \sum_{n=0}^{\infty}  t^n \\
%& \leq A_{\veps}(1-t) + \veps ,
%\end{align*}
%en posant $A_{\veps} = \sum_{n=0}^{N_{\veps}} \abs{ S_n(1) - \varphi(i) }$. Et en se plaçant dans un voisinage (à gauche) de 1 tel que 
%$A_{\veps}(1-t) \leq \veps}$ on a finalement $\abs{\varphi(t) - varphi(1)} \leq 2 \veps$, et ce pour tout $\veps >0$.
%Autrement dit :
%\[ \lim{t \to 1^{-1}} \varphi(t) = varphi(1) , \]
%ce qui est le résultat annoncé. 
%\end{proof}
%
%\begin{coro}
%Soit $\sum a_n z^n$ une série entière à coefficients complexes de rayon de convergence $R$ fini et strictement positif.
%Si la série $\sum a_n R^n$ (resp. $\sum a_n (-1)^n R^n$) est convergente, alors la restriction de la fonction somme $S$ à $]-R, R[$
%se prolonge par continuité en $R$ (resp. en $-R$) et on a :
%\[ S(R) = \sum_{n=0}^{\infty} a_n R^n = \lim_{x \to R} S(x) 
%\text{(resp. } S(-R) = \sum_{n=0}^{\infty} a_n (-1)^, R^n = \lim_{x \to -R} S(x) ) . \]
%\end{coro}

\end{document}